\section*{Цель работы}
Изучение средств и методов измерения параметров электрических цепей (сопротивления, емкости, индуктивности); экспериментальная оценка результатов и погрешностей измерений.

\section*{Задание}
\begin{enumerate}
    \item Ознакомиться с измерительными приборами, представленными на лабораторном стенде.
    \item Выполнить измерение сопротивления резисторов различными типами приборов: измерителем иммитанса, универсальным электронным вольтметром, цифровым вольтметром и комбинированным прибором (тестером).
    \item Провести измерения параметров реактивных элементов (конденсатора и катушки индуктивности) с использованием различных схем замещения.
    \item Оценить погрешности полученных результатов.
\end{enumerate}

\section*{Краткие теоретические сведения}
Результат измерения любого параметра электрической цепи (например, активного сопротивления $R_x$) представляется в виде:
\begin{equation}
    R_x = R_{\text{пр}} \pm \Delta R
\end{equation}
где $R_{\text{пр}}$ --- значение, считанное по шкале или дисплею прибора, $\Delta R$ --- абсолютная погрешность измерения.

Особое внимание следует уделить приборам с неравномерной шкалой (аналоговые омметры). Для них традиционное понятие нормирующего значения, выраженного в единицах измеряемой величины, неприменимо. В качестве нормирующего значения $L_N$ принимают геометрическую длину шкалы в делениях. Предельная абсолютная погрешность в делениях шкалы определяется как:
\begin{equation}
    \Delta L = \frac{k L_N}{100}
\end{equation}
где $k$ --- класс точности прибора. Для перевода этой погрешности в единицы сопротивления используется метод замещения с помощью магазина сопротивлений.

\begin{figure}[hbt]
    \centering
    \includegraphics[width=1\textwidth]{figs/scale_scheme.pdf}
    \caption{Описание шкал}
    \label{fig:scale_scheme.pdf}
\end{figure}

При измерении параметров реактивных компонентов ($C, L$) измерители иммитанса позволяют выбирать параллельную ($p$) или последовательную ($s$) схему замещения. Выбор схемы зависит от добротности $Q$ или тангенса угла потерь $\text{tg} \, \delta$ компонента.

\section*{Инструкция по выполнению работы}

\subsection*{Измерение активного сопротивления}
\begin{enumerate}
    \item Подготовьте приборы к работе согласно инструкциям по эксплуатации.
    \item Для каждого из предложенных резисторов (встроенных в лабораторный модуль) проведите измерения следующими приборами:
    \begin{itemize}
        \item \textit{Измеритель иммитанса:} выберите режим измерения $R$.
        \item \textit{Электронные вольтметры (аналоговый и цифровой):} переключите в режим омметра, выберите подходящий предел измерения.
        \item \textit{Комбинированный прибор (тестер):} установите переключатель в режим $\Omega$.
    \end{itemize}
    \item \textit{Важно:} при работе с аналоговым тестером перед каждым измерением на новом пределе необходимо производить «установку нуля» при замкнутых щупах.
    \item Данные занесите в \cref{tab:resistors_data}.
\end{enumerate}

\subsection*{Оценка погрешности тестера с неравномерной шкалой}
\begin{enumerate}
    \item Подключите магазин сопротивлений к тестеру.
    \item Установите на магазине значение, равное номинальному значению измеряемого резистора.
    \item Отложите по шкале прибора влево и вправо значение $\Delta L$, рассчитанное по формуле (8.2).
    \item С помощью магазина сопротивлений определите значения $R_1$ и $R_2$, соответствующие этим отклонениям.
    \item Зафиксируйте значения в \cref{tab:non_uniform_scale}.
\end{enumerate}

\subsection*{Измерение параметров конденсаторов и катушек}
\begin{enumerate}
    \item Подключите объект измерения к клеммам измерителя иммитанса.
    \item Выполните измерения емкости $C$ и $\text{tg} \, \delta$ для конденсатора, а также индуктивности $L$ и добротности $Q$ для катушки.
    \item Повторите измерения для последовательной ($s$) и параллельной ($p$) схем замещения.
    \item Результаты занесите в \cref{tab:reactive_elements}.
\end{enumerate}

\section*{Инструкция по обработке результатов}
\begin{enumerate}
    \item Рассчитайте абсолютные и относительные погрешности для каждого прибора, используя паспортные данные (классы точности, пределы).
    \item Сравните точность различных методов измерения сопротивления.
    \item Проанализируйте влияние выбора схемы замещения на результаты измерения $C$ и $L$.
\end{enumerate}

% Соответствие изображений:
% Рис. 8.1 из методички соответствует fig:scale_scheme (если бы он был добавлен).